Herein are details of the methods used in Ultra Terminum (UT).
Additionally: details and arguments for UT's novelty, inventiveness, and utility.

UT is a combination of multiple, incremental methods.

The first is \emph{Proof of Reflection}: a method and protocol whereby
a blockchain may utilize an otherwise separate blockchain to
increase its own security and inherit particular security \emph{qualities} of that other blockchain
(which the first blockchain would otherwise never have).
In effect, this is a method for sharing security between blockchains.
If a blockchain uses this method, it does not preclude that blockchain from using it again with another, third blockchain.
Furthermore, this method is universal over all classes of consensus mechanisms that
blockchains may use (e.g., Proof of Work, Proof of Stake, Proof of Authority, etc).

The second is a method for arranging and efficiently operating a set of blockchains using Proof of Reflection to maximize throughput (measured in, e.g., transactions per second), security, and speed.
The result of this method is a construct called \emph{the simplex}.

The third is a method for constructing an unbounded tiling
of simplexes which preserves the security benefits from the second method,
thus providing practically infinite capacity and throughput in a secure environment.

UT claims to solve a well-known problem in the blockchain space, usually called the \emph{scalability trilemma}.
In light of this: the novelty, inventiveness, and utility of UT is somewhat self-evident, as a solution to this problem is highly sought-after.
Additional reasoning and documentation is provided.
